\begin{table}[htbp]
\centering
\small
\begin{threeparttable}
\caption{Spot response of the Taiwan dollar on announcement days}\label{tab:spot}
\begin{tabular}{@{}lp{4.6cm}cccc@{}}
\toprule
& & & \multicolumn{3}{c}{Abnormal change against} \\ \cmidrule(l){4-6}
Window & Event & Raw change & SGD & KRW and SGD & $p$ \\
\midrule
16 Sep 2025 & Industry asks members to develop accounting relief & $-$0.66 & $-$0.50 & $-$0.52 & 0.11 \\
21 Nov 2025 & Industry association proposes the change & +0.64 & +0.72 & +0.63 & 0.07 \\
18 Dec 2025 & Regulator reports to the legislature; public hearing & +0.10 & +0.05 & +0.06 & 0.82 \\
23 Dec 2025 & Draft amendment published & $-$0.35 & $-$0.17 & +0.04 & 0.90 \\
21 Jan 2026 & First insurer resolves to adopt & +0.03 & +0.09 & +0.18 & 0.57 \\
29 Jan 2026 & Amendment issued; large insurers adopt & +0.76 & +0.54 & +0.62 & 0.07 \\
\midrule
\multicolumn{3}{@{}l}{Sum over the six windows} & +0.73 & +1.01 & \\
\multicolumn{3}{@{}l}{Randomisation $p$, one-sided (two-sided)} & 0.33 (0.41) & 0.19 (0.24) & \\
\multicolumn{3}{@{}l}{Standard deviation of two-day abnormal change} & 0.33 & 0.34 & \\
\bottomrule
\end{tabular}
\begin{tablenotes}
\footnotesize
\item Notes: Changes in the log price of the US dollar in local currency, in per cent, over two-day windows beginning on the event date; a positive value is a weaker Taiwan dollar. Abnormal changes subtract the comparison currency's change times the slope of a regression of daily Taiwan-dollar changes on it (slopes 0.47 and 0.27). The $p$-value for each window is the share of the 230 non-event two-day windows with a larger absolute abnormal change (KRW and SGD comparison). Randomisation $p$-values for the sum compare it with 20{,}000 sums of six windows drawn from the non-event windows. Federal Reserve noon buying rates, 2 July 2025 to 30 June 2026 (249 trading days).
\end{tablenotes}
\end{threeparttable}
\end{table}
